% ###################
% ## Filename: SCM_network.tex
% ## Location: D:\Project\Prej_Teaching\ISE\SupplyChainMgt\presentation\scm_07_network
% ## Created: Mar 04, 2011
% ## Last updated: Mar 04, 2011
% ## Figure Location: ..\Figures\
% ## Note: Presentation, Teaching
% ###################


%\documentclass[10pt,presentation]{beamer}
%\documentclass[10pt,handout]{beamer}
\documentclass[10pt,handout,notes=show]{beamer}

%============================================================================================================--
\usepackage{fontspec}
\usepackage{xunicode}
\usepackage{xltxtra}
\usepackage{array}
\usepackage{xcolor}
\usepackage{verbatim}
\usepackage{beamerthemesplit}
\usepackage{beamerbasetemplates}
\usepackage{graphics}
\usepackage{graphicx}
\usepackage{hyperref}
\usepackage{multirow}
\usepackage{tikz,pgfplots}
\usepackage{algorithm}
\usepackage{algorithmic}
\usepackage{soul,xcolor}
\usepackage{colortbl}
\usepackage{multicol}

\XeTeXlinebreaklocale"th"
\defaultfontfeatures{Scale=1.5, Mapping=tex-text} %{Scale=1.23}



\usetheme{Madrid}  % JuanLesPins  Rochester
           %  Berkeley  Palo Alto    with sidebar and top
           % Goettingen  Marburg     with sidebar
           %  Copenhagen  Luebeck  Warsaw

\setbeamertemplate{headline}
{%
    \begin{beamercolorbox}{section in head/foot}
        \vskip-2pt \hskip5pt
        \insertsectionnavigationhorizontal{\textwidth}{}{}
        \hskip5pt \vskip-2pt
    \end{beamercolorbox}
}%
\setbeamertemplate{footline}
{%
    \begin{beamercolorbox}{section in head/foot}
        \vskip-2pt \hskip5pt
        \insertshorttitle \hfill \insertshortdate \hskip10pt \insertframenumber / \inserttotalframenumber
        \vskip-2pt\hskip5pt
    \end{beamercolorbox}
}%
\definecolor{alertedOrange}{RGB}{255,36,0}
\setbeamercolor{alerted text}{fg=alertedOrange}
\useinnertheme{rounded}
\usecolortheme{rose}
%\setbeamercovered{transparent}
\usefonttheme{structuresmallcapsserif}
\definecolor{LRed}{rgb} {1,.8,.8}
\definecolor{LBlue}{rgb}{.8,.8,1}

%%%%%%%%%%%%%%%%%%%%%%%%%%%%%%%%%%%%%%%%%%%%%%%%%%%%%%%%%%%%%%%%%%%%%%%%%%%%%%%
% Notation:
\newtheorem{observation}{Observation}
\newcommand{\figureLocation}[1]{../figures/#1}
\newcommand{\figureLoc}[1]{../figures/#1}
\newcommand{\refLocation}[1]{../../../../Reference/#1}
\newcommand{\NextLine}{\vskip 0.05\textheight}
\newcommand{\NextHalfLine}{\vskip 0.02\textheight}
\newcommand{\RemoveLine}{\vskip -0.1\textheight}
\newcommand{\RemoveHalfLine}{\vskip -0.02\textheight}
\newcommand{\BackTab}{\hskip -0.15\textwidth}
\newcommand{\HalfBackTab}{\hskip -0.11\textwidth}
\newcommand{\DualFigureSize}{1.10\textwidth}
\newcommand{\MiniPageSize}{0.52\textwidth}
\newcommand{\code}[1]{\textcolor{blue}{#1} }
\newcommand{\rightNote}[2]{\hfill \uncover<#1>{\textbf{\code{\#\# {#2}  }}}}
\newcolumntype{P}[1]{>{\raggedright\arraybackslash}p{#1}<{}}

\author[Kittithreerapronchai]{%
  Oran~Kittithreerapronchai\inst{1} }
\institute[IE Chula]{
  \inst{1}%
  Department of Industrial Engineering, Chulalongkorn University\\
  Bangkok 10330 THAILAND \\
  }
\subject{}
\date{}

\setbeamertemplate{enumerate items}[default]
\setbeamertemplate{bibliography item}[text]
\bibliographystyle{alpha}
\newcommand{\textFont}[1]{\footnotesize{#1}}
\setcounter{MaxMatrixCols}{20}

%%%--------------------------------

\makeatletter
\defbeamertemplate{note page}{infolines}
{%
  {%
    \small
    \insertvrule{.25\paperheight}{white!90!black}
    \vskip-.25\paperheight
    \nointerlineskip
    \vbox{
      \hfill\insertslideintonotes{0.25}\hskip-\Gm@rmargin\hskip0pt%
      \vskip-0.25\paperheight%
      \nointerlineskip
      \begin{pgfpicture}{0cm}{0cm}{0cm}{0cm}
        \begin{pgflowlevelscope}{\pgftransformrotate{90}}
          {\pgftransformshift{\pgfpoint{-2cm}{0.2cm}}%
          \pgftext[base,left]{\footnotesize\the\year-\ifnum\month<10\relax0\fi\the\month-\ifnum\day<10\relax0\fi\the\day}}
        \end{pgflowlevelscope}
      \end{pgfpicture}}
    \nointerlineskip
    \vbox to .25\paperheight{\vskip0.5em
      \hbox{\insertshorttitle[width=8cm]}%
      \setbox\beamer@tempbox=\hbox{\insertsection}%
      \hbox{\ifdim\wd\beamer@tempbox>1pt{\hskip4pt\raise3pt\hbox{\vrule
            width0.4pt height7pt\vrule width 9pt
            height0.4pt}}\hskip1pt\hbox{\begin{minipage}[t]{7.5cm}\def\breakhere{}\insertsection\end{minipage}}\fi%
      }%
      \setbox\beamer@tempbox=\hbox{\insertsubsection}%
      \hbox{\ifdim\wd\beamer@tempbox>1pt{\hskip17.4pt\raise3pt\hbox{\vrule
            width0.4pt height7pt\vrule width 9pt
            height0.4pt}}\hskip1pt\hbox{\begin{minipage}[t]{7.5cm}\def\breakhere{}\insertsubsection\end{minipage}}\fi%
      }%
      \setbox\beamer@tempbox=\hbox{\insertshortframetitle}%
      \hbox{\ifdim\wd\beamer@tempbox>1pt{\hskip30.8pt\raise3pt\hbox{\vrule
            width0.4pt height7pt\vrule width 9pt
            height0.4pt}}\hskip1pt\hbox{\insertshortframetitle[width=7cm]}\fi%
      }%
      \vfil}%
  }%
  \vskip.25em
  \nointerlineskip
  %\setmainfont{TH Sarabun New}
  \begin{minipage}{\textwidth} % this is an addition
  \footnotesize{\fontspec{TH Sarabun New} \insertnote}
  \end{minipage}               % this is an addition
  %\setmainfont{Computer Modern}

}
\makeatother

\setbeamertemplate{note page}[infolines]

\title[LSCM v3.0]{Lecture 08:  \\ Network Design and Site Selection \\ ~~ }
\hypersetup{
  pdftitle={09 Network Design and Site Selection},
  pdfsubject={2104509lscm & 2184408scm handout},
  pdfproducer={XeLatex/WinEdt},
  pdfcreator={oran.k@chula.ac.th}
}

\begin{document}

%%==============================================================================================================
\begin{frame}[t]
  \titlepage \label{Slide:Title}

  \NextLine   \NextLine
  \begin{flushright}
    \tiny{ \textbf{last updated:} \today } \hskip3ex
  \end{flushright}

  \note{
    \begin{itemize}
      \item introduce to network model $\rightarrow$ real SCM job
      \item business logistic $\rightarrow$ know practical logistics
      \item SCM $\rightarrow$ use model to predict/design
    \end{itemize}
  }
\end{frame}



%%==============================================================================================================
\begin{frame}[t]
  \frametitle{Outline}

  \tableofcontents
  \vfill
  \begin{flushright}
    \tiny{ \textbf{Key Ref.:}~\cite{jacobs2010operations}~\cite{ballou2007evolution} \cite{chopra2007supply} \cite{goetschalckx2011supply}} \hskip3ex
  \end{flushright}
  \note{
    \begin{itemize}
      \item important and motivation of LSCM~\cite{jacobs2010operations}~\cite{ballou2007evolution} \cite{chopra2007supply} \cite{goetschalckx2011supply}
      \item LSCM is an elective class
    \end{itemize}
  }
\end{frame}

%%=================================================================================================
% Standardize how tikz displays vendors, crossdock, and stores
\usetikzlibrary{shapes}
\tikzstyle{vendor}=[rectangle,draw=black!80,thick,inner sep=0pt,minimum size=5mm]
\tikzstyle{store}=[circle,draw=blue!80,fill=blue!60,thick,inner sep=0pt,minimum size=5mm]
\tikzstyle{xdock}=[diamond,draw=red!80,fill=red!60,thick,inner sep=0pt,minimum size=5mm]
\tikzstyle{tl}=[->,thick, black]
\tikzstyle{ltl}=[->,thick, gray, dotted]
\tikzstyle{tl2way}=[<->,thick, black]
\tikzstyle{ltl2way}=[<->,thick, gray, dotted]
\tikzstyle{aLeft}=[above left=0.4cm]
\tikzstyle{aRight}=[above right=0.4cm]
\tikzstyle{bLeft}=[below left=0.4cm]
\tikzstyle{bRight}=[below right=0.4cm]

%%=================================================================================================
\section[Network 101]{Introduction to Distribution Network}
%%=================================================================================================
\begin{frame}

  \frametitle{Why do we transportation?}
  \begin{itemize}
    \item \textbf{Economics:} scales, low cost, JIT
    \item \textbf{Geographic} raw material, markets, resources (labor)
    \item \textbf{Proprietary:} special process, softwares, R\&D
    \item \textbf{Environment:} danger, hazard
  \end{itemize}

  \NextLine
  \begin{alertblock}{\small{Trade-Off in transportation design}}{}
    \begin{center}
    \begin{itemize}
      \item Inventory Cost \textbf{VS} Transportation Cost
      \item Responsive \textbf{VS} Transportation Cost
    \end{itemize}
    \end{center}
  \end{alertblock}

  \note{
      \begin{itemize}
        \item \textbf{Responsibility:} WalMart on illegal immigrant, sexism, imbalance trade
      \end{itemize}
 }
 \end{frame}


%=================================================================================================
\begin{frame}
  \frametitle{What should you know about network?}

  \begin{itemize}
    \item \textbf{Goal:} \alert<2->{minimizing costs}, min-max service levels, min-max hazard/risk
    \item \textbf{Echelon Level:} warehouse--DC, retailer--customer
    \item \textbf{Infrastructure:} \alert<2->{transit facility}, inventory, trucks, equipments, WMS
    \item \textbf{Routing:} dynamic \textbf{VS} static
    \item \textbf{Visiting Policy:} single destination \textbf{VS} \alert<2->{milk-run}
  \end{itemize}

  \NextLine
  \begin{alertblock}{}{}
    \textbf{Distribution network}: steps taken to \alert{move} and \alert{store} products from suppliers to targeted customers in a supply chain
  \end{alertblock}

  \note{
      \begin{itemize}
        \item Distribution directly affects cost and the customer experience and therefore drives profitability
        \item must aligns with strategy
        \item information
      \end{itemize}
 }
 \end{frame}

%=================================================================================================
\begin{frame}
  \frametitle{Single Destination}

  \begin{center}
    \begin{tikzpicture}[xscale=4,yscale=2]

      \node[store] (v1) at (0,1)  {};
      \node[store] (vn) at (0,0) {};
      \node[store] (st1) at (1,1) {};
      \node[store] (st2) at (1,.5) {};
      \node[store] (st3) at (0.5,.0) {};
      \node[store] (stn) at (1,0) {};
      \node[xdock] (xd1) at (0.5,0.5) {};

      \draw [tl] (xd1) -- (v1);
      \draw [tl] (xd1) -- (st1);
      \draw [tl] (xd1) -- (st2);
      %\draw [->,thick, red, dotted]  (st2) -- (xd1) ;
      \draw [tl] (xd1) -- (stn);
      \draw [tl] (xd1) -- (st3);
      \draw [tl] (xd1) -- (vn);
      %\draw [->,thick, red, dotted]  (vn) -- (xd1) ;
    \end{tikzpicture}
  \end{center}

  \NextHalfLine
  \begin{itemize}
    \item<1-> \textbf{Idea:} visiting \alert{single location} in each trip
    \item<2-> \textbf{Pros:} simple, high-level model
    \item<2-> \textbf{Cons:} \alert{fleet management}, \# trucks
    \item<2-> \textbf{Examples:} \emph{rarely seen in pure form} why?
  \end{itemize}

  \note{
      \begin{itemize}
        \item A company focuses on inbound b/c easy to pressure supplier than customers
        \item A company must find customer to return trip
      \end{itemize}
 }
 \end{frame}

%=================================================================================================
\begin{frame}
  \frametitle{Milk Run Network}

  \begin{center}
    \begin{tikzpicture}[xscale=4,yscale=2]

      \node[store] (v1) at (0,1)  {};
      \node[store] (vn) at (0,0) {};
      \node[store] (st1) at (1,1) {};
      \node[store] (st2) at (1,.5) {};
      \node[store] (st3) at (0.5,.0) {};
      \node[store] (stn) at (1,0) {};
      \node[xdock] (xd1) at (0.5,0.5) {};

      \draw [tl] (xd1) -- (v1);
      \draw [tl] (v1) -- (st1);
      \draw [tl] (st1) -- (st2);
      \draw [->,thick, red, dotted]  (st2) -- (xd1) ;
      \draw [tl] (xd1) -- (stn);
      \draw [tl] (stn) -- (st3);
      \draw [tl] (st3) -- (vn);
      \draw [->,thick, red, dotted]  (vn) -- (xd1) ;
    \end{tikzpicture}
  \end{center}

  \NextHalfLine
  \begin{itemize}
    \item<1-> \textbf{Idea:} visiting \alert{many locations} in one trip
    \item<2-> \textbf{Pros:} high utilization, low inventory
    \item<2-> \textbf{Cons:} \alert{high shipment frequency}, required \alert{planning/routing}
    \item<2-> \textbf{Examples:} 7-Eleven (DC to stores)
  \end{itemize}

  \note{
      \begin{itemize}
        \item A company focuses on inbound b/c easy to pressure supplier than customers
      \end{itemize}
 }
 \end{frame}

%=================================================================================================
\begin{frame}
  \frametitle{Direct Shipment Network}

  \begin{center}
    \begin{tikzpicture}[xscale=4,yscale=2]

      \node[vendor] (v1) at (0,1)  {};
      \node[vendor] (vn) at (0,0) {};
      \node[store] (st1) at (1,1) {};
      \node[store] (stn) at (1,0) {};

      \node at (0,0.5) {$\vdots$};
      \node at (0,-0.5) {Origins};

      \node at (1,0.5) {$\vdots$};
      \node at (1,-0.5) {Destinations};

      \draw [ltl] (v1) -- (st1);
      \draw [ltl] (v1) -- (stn);
      \draw [ltl] (vn) -- (st1);
      \draw [ltl] (vn) -- (stn);

    \end{tikzpicture}
  \end{center}

  \NextHalfLine
  \begin{itemize}
    \item<2-> \textbf{Pros:} simplest network, low transit time, required \alert{no facility} or IT
    \item<2-> \textbf{Cons:} more vehicles, no consolidation, \alert{low utilization}
    \item<2-> \textbf{Examples:} urgent delivery, full truckload shipment
  \end{itemize}

  \note{
      \begin{itemize}
        \item single shipment to the destination and done
        \item can combine with milk-run
        \item it should be nearly full truckload
        \item Retailer carries no inventory.
        \item eBags (no inventory at all), nordstrom.com, WW Grainger (use this model for slow-moving products)
        \item Benefits: mfg gets to aggregate demands and inventory. -> high product availability with low inventory
        \item Ownership structure of inventory at the mfg -> better if inventory not allocated to particular customers
        \item Benefits are high for high-value, low-demand, unpredictable demand.
      \end{itemize}
 }
 \end{frame}

%=================================================================================================

\begin{frame}
  \frametitle{Hub-and-Spoke Network}

  \begin{center}
    \begin{tikzpicture}[xscale=4,yscale=2]

      \node[vendor] (v1) at (0,1)  {};
      \node[vendor] (vn) at (0,0) {};
      \node[store] (st1) at (1,1) {};
      \node[store] (stn) at (1,0) {};
      \node[xdock] (xd1) at (0.25,0.5) {};
      \node[xdock] (xd2) at (0.75,0.5) {};

      \node at (0,0.5) {$\vdots$};
      \node at (0,-0.5) {Origins};

      \node at (1,0.5) {$\vdots$};
      \node at (1,-0.5) {Destinations};

      \node at (0.5,-0.5) {Hubs};

      \draw [ltl] (v1) -- (xd1);
      \draw [ltl] (vn) -- (xd1);
      \draw [ltl] (xd2) -- (st1);
      \draw [ltl] (xd2) -- (stn);
      \draw[tl] (xd1) -- (xd2);
    \end{tikzpicture}
  \end{center}

    \NextHalfLine
  \begin{itemize}
    \item<1-> \textbf{Idea:} hub as a temporary destination/origin to consolidate
    \item<2-> \textbf{Pros:} \alert{economies of scale} $\rightarrow$ high utilization
    \item<2-> \textbf{Cons:} required \alert{facilities} and sorting machine, \alert{higher transit time}
    \item<2-> \textbf{Examples:} Airlines, Mails, 3PL Companies
  \end{itemize}

  \note{
      \begin{itemize}
        \item A company focuses on inbound b/c easy to pressure supplier than customers
        \item what if mail in one location
      \end{itemize}
 }
 \end{frame}

%=================================================================================================
\begin{frame}
  \frametitle{Single-Hub Network: Crossdocking}

  \begin{center}
    \begin{tikzpicture}[xscale=4,yscale=2]

      \node[vendor] (v1) at (0,1)  {};
      \node[vendor] (vn) at (0,0) {};
      \node[store] (st1) at (1,1) {};
      \node[store] (stn) at (1,0) {};
      \node[xdock] (xd1) at (0.5,0.5) {};

      \node at (0,0.5) {$\vdots$};
      \node at (0,-0.5) {Origins};

      \node at (1,0.5) {$\vdots$};
      \node at (1,-0.5) {Destinations};

      \node at (0.5,-0.5) {Hubs};

      %\draw [ltl](v1) -- (st1);
      %\draw [ltl](vn) -- (stn);
      \draw [tl] (v1) -- (xd1);
      \draw [tl] (vn) -- (xd1);
      \draw [tl] (xd1) -- (st1);
      \draw [tl] (xd1) -- (stn);
    \end{tikzpicture}
  \end{center}

  \NextHalfLine
  \begin{itemize}
    \item<1-> \textbf{Idea:} \alert{variation} of hub-and-spoke for relatively high flow
    \item<2-> \textbf{Pros:} low \alert{transit time}, economy of scale, high utilization
    \item<2-> \textbf{Cons:} required facilities and sorting
    \item<2-> \textbf{Examples:} Crossdocking (Wal-Mart, Home depot), large 3PL Companies
  \end{itemize}

  \note{
      \begin{itemize}
        \item A company focuses on inbound b/c easy to pressure supplier than customers
      \end{itemize}
 }
 \end{frame}


%=================================================================================================
\begin{frame}
  \frametitle{Why do we need a model?}
  \begin{itemize}
    \item understand and analyze network
    \item predict and determine decisions
  \end{itemize}

  \NextHalfLine
  \textsc{Quantitative Model Components}
  \begin{itemize}
    \item \textbf{Assumptions:} \alert{simplify factors}, capture important parts
    \item \textbf{Data:} verify parameters and distributions
    \item \textbf{Model:} describe key \alert{relationships} and interactions
    \item \textbf{Tools:} spreadsheet, optimization, simulation
    \item \textbf{Conclusion:} suggestions and \alert{insights}, not numerical value
  \end{itemize}

  \NextHalfLine
  \textsc{Network optimization model:} \uncover<3->{\alert{graph theory} based model}
  \begin{itemize}
      \item<2-> Shortest Path Problem; Transportation Problem
      \item<2-> Capacitated Plant Location Problem
      \item<2-> Traveling Salesman Problem (TSP)
  \end{itemize}


  \note{
      \begin{itemize}
        \item \textbf{so far:} introduction on terminology
        \item \textbf{why model?:} quantify idea before implement
        \item \textbf{assumption:} too simplify is trivial; too complicate is uncomprehensible
        \item \textbf{data:} data can easily changed so sensitivity analysis
      \end{itemize}
   }
 \end{frame}

%=================================================================================================
\section[Review]{Graph Theory Review and Transportation Model}
%=================================================================================================
% Standardize how tikz displays vendors, crossdock, and stores
\usetikzlibrary{shapes,arrows,positioning}
\tikzstyle{node}=[circle,draw=black!80,thick,fill=yellow!15,inner sep=0pt,minimum size=5mm]
\tikzstyle{dNode}=[circle,draw=black!40,fill=blue!15,inner sep=0pt,minimum size=8mm]
\tikzstyle{eEdge}=[->,black!70,dotted,thick,>=triangle 45]
\tikzstyle{aEdge}=[->, black,thick,>=triangle 45]
\tikzstyle{neEdge}=[-,black!70,dotted,thick,>=triangle 45]
\tikzstyle{naEdge}=[-, black,thick,>=triangle 45]
\tikzstyle{edgeText}=[pos=0.45,auto, blue,sloped]
\newcommand{\tikztext}[1]{\footnotesize{#1}}
%===================================================================================================
\begin{frame}
  \frametitle{Introduction to Graph Theory}

  \RemoveLine \RemoveHalfLine
  \begin{center}
    \begin{tikzpicture}[xscale=2,yscale=2]
      \node[node] (n1) at (0.0,1.5) {\tikztext{1}};
      \node[node] (n2) at (0.0,0.0) {\tikztext{2}};
      \node[node] (n3) at (1.5,0.0) {\tikztext{3}};
      \node[node] (n4) at (1.5,1.5) {\tikztext{4}};

      \draw[eEdge] (n1) -- (n2) ;
      \draw[eEdge] (n2) -- (n3) ;
      \draw[eEdge] (n3) -- (n4) ;
      \draw[eEdge] (n4) -- (n1) ;
      \path[eEdge] (n4) edge[bend left] (n3)    ;
    \end{tikzpicture}
  \end{center}


  \begin{itemize}
    \item \textbf{Node:} (vertex) presentation of stage
    \item<2->[] \hskip5em $\mathcal{N} \in\ \{1,2,3,4 \}$
    \item \textbf{Edge:} (arc) connecting line between nodes
    \item<2->[] \hskip5em $\mathcal{E} \in\ \{ (1,2),(2,3),(3,4),(4,1),(4,3) \}$
  \end{itemize}

  \note{
    \begin{itemize}
      \item a
    \end{itemize}
  }
\end{frame}

%=================================================================================================
\begin{frame}
  \frametitle{Parameters in Graph}

  \RemoveLine \RemoveHalfLine
  \begin{center}
    \begin{tikzpicture}[xscale=2,yscale=2]
      \node[node,label=above:\tiny{10}] (n1) at (0.0,1.5) {\tikztext{1}};
      \node[node,label=below:\tiny{20}] (n2) at (0.0,0.0) {\tikztext{2}};
      \node[node,label=below:\tiny{-15}] (n3) at (1.5,0.0) {\tikztext{3}};
      \node[node,label=above:\tiny{-15}] (n4) at (1.5,1.5) {\tikztext{4}};

      \draw[eEdge] (n1) -- node[edgeText,above] {\tiny{$\infty,\$5$}}  (n2) ;
      \draw[eEdge] (n2) -- node[edgeText] {\tiny{$10,\$1$}}  (n3) ;
      \draw[eEdge] (n3) -- node[edgeText,below] {\tiny{$30,\$2$}}  (n4) ;
      \draw[eEdge] (n4) -- node[edgeText] {\tiny{$15,\$2$}}  (n1) ;
      \path[eEdge] (n4) edge[bend left] node[edgeText] {\tiny{$8,\$4$}} (n3)  ;
    \end{tikzpicture}
  \end{center}

  \begin{itemize}
    \item \textbf{Demand/Suppy:} demands and supply at each node, e.g. demands at node 1 is 10
    \item \textbf{Cost:} incurred cost per a unit of flow, e.g., cost at edge $(1,2)$ is 5
    \item \textbf{Capacity:} unit flow allow at each edge, e.g., flow at edge $(1,2)$ must be less than $\infty$
  \end{itemize}

  \note{
    \begin{itemize}
      \item a
    \end{itemize}
  }
\end{frame}

%=================================================================================================
\begin{frame}
  \frametitle{Simple Network }

  \small{
   A network with 6 vertices, $\mathcal{N} \in\ \{1,2,3,4,5,6 \}$, and $\{-12,24,-14,-16,30,-12\}$ demands/suppies at vertices and the following cost and capacity at each edge. Construct the graph of this network
  \begin{center}
  \begin{tabular}{lrr}\hline
    Edge    & Cost  & Max Capacitiy \\ \hline
    (1,2)   &  \$2    &  6 \\
    (2,4)   &  \$2    &  6 \\
    (4,6)   &  \$2    &  6 \\
    (5,6)   &  \$2    &  6 \\
    (3,5)   &  \$2    &  6 \\
    (1,3)   &  \$2    &  6 \\
    (1,5)   &  \$5    &  $\infty$ \\
    (2,3)   &  \$1    &  2 \\
    (2,6)   &  \$5    &  $\infty$ \\
    (4,5)   &  \$1    &  2 \\ \hline
  \end{tabular}
  \end{center}
  }

  \note{
    \begin{itemize}
      \item a
    \end{itemize}
  }
\end{frame}

%=================================================================================================
\begin{frame}
  \frametitle{Graph Representation  }

  \RemoveLine \RemoveHalfLine
  \begin{center}
    \begin{tikzpicture}[xscale=2,yscale=2]
      \node[node,label=above:\tiny{-12}]  (n1) at (0.0,1.4) {\tikztext{1}};
      \node[node,label=above:\tiny{24}] (n2) at (0.8,2.0) {\tikztext{2}};
      \node[node,label=below:\tiny{-14}]  (n3) at (0.8,0.8) {\tikztext{3}};
      \node[node,label=above:\tiny{-16}]  (n4) at (2.2,2.0) {\tikztext{4}};
      \node[node,label=below:\tiny{30}] (n5) at (2.2,0.8) {\tikztext{5}};
      \node[node,label=below:\tiny{-12}] (n6) at (3.0,1.4) {\tikztext{6}};

      \only<2|handout:0>{
      \draw[eEdge] (n1) -- node[edgeText] {\tiny{$6,\$2$}}  (n2) ;
      \draw[eEdge] (n2) -- node[edgeText] {\tiny{$6,\$2$}}  (n4) ;
      \draw[eEdge] (n4) -- node[edgeText] {\tiny{$6,\$2$}}  (n6) ;
      \draw[eEdge] (n5) -- node[edgeText] {\tiny{$6,\$2$}}  (n6) ;
      \draw[eEdge] (n3) -- node[edgeText] {\tiny{$6,\$2$}}  (n5) ;
      \draw[eEdge] (n1) -- node[edgeText] {\tiny{$6,\$2$}}  (n3) ;
      \draw[eEdge] (n1) -- node[edgeText] {\tiny{$\infty,\$5$}}  (n5) ;
      \draw[eEdge] (n2) -- node[edgeText] {\tiny{$\infty,\$5$}}  (n6) ;
      \draw[eEdge] (n2) -- node[edgeText,above] {\tiny{$2,\$1$}}  (n3) ;
      \draw[eEdge] (n4) -- node[edgeText,above] {\tiny{$2,\$1$}}  (n5) ;
      }

    \end{tikzpicture}
  \end{center}

  \note{
    \begin{itemize}
      \item \textbf{optimal solution:} N/A
    \end{itemize}
  }
\end{frame}

%=================================================================================================
\begin{frame}
  \frametitle{Shortest Path  }

  \begin{itemize}
    \item \textbf{What:} model describes the shortest/fastest \alert{path} to transport product(s) from origin to destination
    \item \textbf{Idea:} minimizing total times/distances
    \item \textbf{Assumption:} single product, no capacity
    \item \textbf{Components:}
    \begin{itemize}
      \item single origin/ single destination
      \item transit nodes
    \end{itemize}
    \item \textbf{Application:} GPS
  \end{itemize}

  \note{
    \begin{itemize}
      \item what is the optimal course of action to fulfil demand
    \end{itemize}
  }
\end{frame}
%=================================================================================================
\begin{frame}
  \frametitle{Shortest Path }

\small{ [Winston 2003. Example 8.1] Based on the following network, find the shortest way to transport products from Plant to City.
  }
  \RemoveLine
  \begin{center}
    \begin{tikzpicture}[xscale=2,yscale=2]
      \node[node,label=left:\tiny{Plant}] (n1) at (0.0,1.4) {\tikztext{1}};
      \node[node] (n2) at (0.8,2.0) {\tikztext{2}};
      \node[node] (n3) at (0.8,0.8) {\tikztext{3}};
      \node[node] (n4) at (2.2,2.0) {\tikztext{4}};
      \node[node] (n5) at (2.2,0.8) {\tikztext{5}};
      \node[node,label=right:\tiny{City}] (n6) at (3.0,1.4) {\tikztext{6}};

    \only<1|handout:1>{
      \draw[eEdge] (n1) -- node[edgeText] {\tiny{$\$4$}}  (n2) ;
      \draw[eEdge] (n2) -- node[edgeText] {\tiny{$\$3$}}  (n4) ;
      \draw[eEdge] (n4) -- node[edgeText] {\tiny{$\$2$}}  (n6) ;
      \draw[eEdge] (n5) -- node[edgeText] {\tiny{$\$2$}}  (n6) ;
      \draw[eEdge] (n3) -- node[edgeText] {\tiny{$\$3$}}  (n5) ;
      \draw[eEdge] (n1) -- node[edgeText] {\tiny{$\$3$}}  (n3) ;
      \draw[eEdge] (n2) -- node[edgeText] {\tiny{$\$2$}}  (n5) ;
    }

    \only<2-|handout:0>{
      \draw[aEdge] (n1) -- node[edgeText] {}  (n2) ;
      \draw[aEdge] (n5) -- node[edgeText] {}  (n6) ;
      \draw[aEdge] (n3) -- node[edgeText] {}  (n5) ;
      \draw[aEdge] (n1) -- node[edgeText] {}  (n3) ;
      \draw[aEdge] (n2) -- node[edgeText] {}  (n5) ;
    }
    \end{tikzpicture}
  \end{center}

    \begin{itemize}
      \item \textbf{optimal:} \$8
      \item \textbf{path 1:} 1 $\rightarrow$ 2 $\rightarrow$ 5 $\rightarrow$ 6
      \item \textbf{path 2:} 1 $\rightarrow$ 3 $\rightarrow$ 5 $\rightarrow$ 6
    \end{itemize}


  \note{
    \begin{itemize}
      \item what is the optimal course of action to fulfil demand
    \end{itemize}
  }
\end{frame}

%=================================================================================================
\begin{frame}
  \frametitle{Shortest Path  }

  \begin{flushleft} \small{
    \begin{alignat*}{3}
        \min\ z =   & 4x_{1,2} + 3x_{1,3}+  3x_{2,4} +  2x_{2,5} + 3x_{3,5} + 2x_{4,6} + 2x_{5,6} \\
      \end{alignat*}
      \begin{alignat*}{3}
        \mbox{Node 1;~~~}      & -x_{1,2} - x_{1,3}             &=\ &-1 \\
        \mbox{Node 2;~~~}      & x_{1,2} -  x_{2,4} -  x_{2,5}  &=\ &0 \\
        \mbox{Node 3;~~~}      & x_{1,3} + x_{3,5}              &=\ &0 \\
        \mbox{Node 4;~~~}      & x_{2,4} - x_{4,6}              &=\ &0 \\
        \mbox{Node 5;~~~}      & x_{2,5} + x_{3,5} - x_{5,6}    &=\ &0 \\
        \mbox{Node 6;~~~}      & x_{4,6} + x_{5,6}              &=\ &1 \\
        \mbox{Non Neg;~~~}     & x_{i,j} &\geq\ & 0 \qquad  \forall\ (i,j) \in\ \mathcal{E}
      \end{alignat*}
  }
  \end{flushleft}

  \textbf{Observations}
  \begin{itemize}
    \item supply = demand; $\sum$ RHS = 0
    \item each DV appears twice in model
  \end{itemize}

  \note{
    \begin{itemize}
      \item what is the optimal course of action to fulfil demand
    \end{itemize}
  }
\end{frame}

%=================================================================================================
\begin{frame}
  \frametitle{Transportation Problem  }

  \begin{itemize}
    \item \textbf{What:} model describes the cheapest way to  \alert{allocate} products
    \item \textbf{Idea:} minimizing total costs while satisfying all demands and supplies
    \item \textbf{Assumption:} \alert{single product}, \alert{assignment} (no routing nor capacity)
    \item \textbf{Components:}
    \begin{itemize}
        \item two groups of facilities/agents
        \item limited supplies and demands
        \item limited transportation capacity \NextLine
    \end{itemize}
    \item \textbf{Application:} Parcel shipment
    \item \textbf{Other Names} Demand Allocation
  \end{itemize}


  \note{
    \begin{itemize}
      \item what is the optimal course of action to fulfil demand
    \end{itemize}
  }
\end{frame}

%=================================================================================================
\begin{frame}
  \frametitle{Transportation Problem  }

  \small{ [Winston 2003. Question 7.1.1] A company supplies goods to three customers. Each of which requires the goods up to 30 units. The company has two warehouses. Warehouse 1 has 40 units available and warehouse 2 has 30 unit available the cost of shipping one unit from warehouse to customer are listed as follows:
  \begin{center}
  \begin{tabular}{lrrr} \hline
        ~ & \multicolumn{3}{c}{To} \\ \cline{2-4}
    From        & Customer 1  & Customer 2  & Customer 3 \\ \hline
    Warehouse 1 & \$15 & \$35 & \$25 \\
    Warehouse 2 & \$10 & \$50 & \$40 \\ \hline
  \end{tabular}
  \end{center}
  Construct the graph representative of this network and find the optimal allocation
  }

  \note{
    \begin{itemize}
      \item
    \end{itemize}
  }
\end{frame}

%=================================================================================================
\begin{frame}
  \frametitle{Transportation Problem: Graph }

  \RemoveLine \RemoveHalfLine
  \begin{center}
    \begin{tikzpicture}[xscale=2,yscale=2]
      \node[node,label=left:\tiny{-40}]  (w1) at (0.0,2.0) {\tikztext{1}};
      \node[node,label=left:\tiny{-30}] (w2) at (0.0,1.0) {\tikztext{2}};
      \uncover<2-|handout:0>{\node[node,label=left:\tiny{-20}] (w0) at (0.0,0.0) {\tikztext{\tiny{dummy}}}; }
      \node[node,label=right:\tiny{30}]  (c1) at (3.0,2.0) {\tikztext{1}};
      \node[node,label=right:\tiny{30}] (c2) at (3.0,1.0) {\tikztext{2}};
      \node[node,label=right:\tiny{30}] (c3) at (3.0,0.0) {\tikztext{3}};

      \node (pName) at (0.0,-0.75) {\tikztext{Warehouse}};
      \node (cName) at (3.0,-0.75) {\tikztext{Customer}};

    \uncover<3-|handout:0>{
      \draw[eEdge] (w1) -- node[edgeText] {\tiny{$\$15$}}  (c1) ;
      \draw[eEdge] (w1) -- node[edgeText] {\tiny{$\$35$}}  (c2) ;
      \draw[eEdge] (w1) -- node[edgeText] {\tiny{$\$25$}}  (c3) ;
      \draw[eEdge] (w2) -- node[edgeText] {\tiny{$\$10$}}  (c1) ;
      \draw[eEdge] (w2) -- node[edgeText] {\tiny{$\$50$}}  (c2) ;
      \draw[eEdge] (w2) -- node[edgeText] {\tiny{$\$40$}}  (c3) ;
      \draw[eEdge] (w0) -- node[edgeText] {\tiny{}}  (c1) ;
      \draw[eEdge] (w0) -- node[edgeText] {\tiny{}}  (c2) ;
      \draw[eEdge] (w0) -- node[edgeText] {\tiny{}}  (c3) ;
    }
    \end{tikzpicture}
  \end{center}

  \note{
    \begin{itemize}
      \item a
    \end{itemize}
  }
\end{frame}
%=================================================================================================
\begin{frame}
  \frametitle{Transportation Problem: Model }

  \begin{flushleft}
  \small{
    \begin{alignat*}{2}
    \min\ z     & =& 15x_{w1,c1}+35x_{w1,c2}+25x_{w1,c3} +10x_{w2,c1}+50x_{w2,c3}+40x_{w2,c3}
    \end{alignat*}
    \begin{alignat*}{3}
    \mbox{s.t.} & & & \\
    \mbox{node w1;~~} &  -x_{w1,c1} - x_{w1,c2} - x_{w1,c3} ~&=\ & -40\\
    \mbox{node w2;~~} &  -x_{w2,c1} - x_{w2,c2} - x_{w2,c3} ~&=\ & -30\\
    \mbox{node 0;~~} &  -x_{0,c1} - x_{0,c2} - x_{0,c3} &=\ & -20\\
    \mbox{node c1;~~} &  x_{w1,c1} + x_{w2,c1}+x_{0,c1}  &=\ & 30 \\
    \mbox{node c2;~~} &  x_{w1,c2} + x_{w2,c2}+x_{0,c2}  &=\ & 30 \\
    \mbox{node c3;~~} &  x_{w1,c3} + x_{w2,c3}+x_{0,c3}  &=\ & 30 \\
    \mbox{non-neg.;~~}&  x_{i,j}                            &\geq\ & 0 \qquad \forall\ (i,j) \in\ \mathcal{E}
    \end{alignat*}
  }
  \end{flushleft}

  \note{
    \begin{itemize}
      \item a
    \end{itemize}
  }
\end{frame}

%=================================================================================================
\begin{frame}
  \frametitle{Capacitated Plant Location }

  \begin{itemize}
    \item \textbf{What:} model describes the cheapest way to \alert{setup facilities} and \alert{their allocations}
    \item \textbf{Idea:} minimizing total costs while satisfying all demands and supplies/capacities
    \item \textbf{Assumption:} \alert{single product}, \alert{assignment} (no routing)
    \item \textbf{Components:}
    \begin{itemize}
        \item set of possible location
        \item limited supplies and demands
        \item limited transportation capacity
        \item unable to supply, unless facility is built \NextLine
    \end{itemize}
    \item \textbf{Application:} Preliminary network design, Parcel-based network
  \end{itemize}

  \note{
    \begin{itemize}
      \item a
    \end{itemize}
  }

\end{frame}

%=================================================================================================
\begin{frame}
  \frametitle{Example of CPL  Problem }

  \footnotesize{[\emph{Chopra \& Meindl 2010.}] SunOil, a manufacturer of petrochemical production, wants to set up facilities to supply demand in each region. The facility, production, transportation costs are listed as follows:
  \begin{center}
  \begin{tabular}{lrrrrrclrlr} \hline \hline
       & \multicolumn{5}{c}{prod. \& trans. costs} & ~& \multicolumn{2}{c}{low cap.} & \multicolumn{2}{c}{high cap.} \\ \cline{2-6} \cline{8-9} \cline{10-11}
    Region & N. Amer & S. Amer & Europe & Asia & Africa & & fixed & cap  & fixed & cap \\ \hline
    N. Amer      & 81 & 92 & 101 & 130 & 115 & & 6,000 & 10 &  9,000 & 20 \\
    S. Amer      & 117& 77 & 108 & 98  & 100 & & 4,500 & 10 &  6,750 & 20 \\
    Europe        & 102& 105& 95  & 119 & 111 & & 6,500 & 10 &  9,750 & 20 \\
    Asia          & 115& 125& 90  & 59  & 74  & & 4,100 & 10 &  6,150 & 20 \\
    Africa        & 142& 100& 103 & 105 & 71  & & 4,000 & 10 &  6,000 & 20 \\
    \\ \cline{1-6}
    Demand        & 12 & 8  & 14  & 16  & 7 \\ \hline \hline
  \end{tabular}
  \end{center}
  Find the optimal facility location of SunOil.
  }

  \note{
    \begin{itemize}
      \item what is the optimal course of action to fulfil demand
    \end{itemize}
  }
\end{frame}


%=================================================================================================
\begin{frame}
  \frametitle{Graph of CPL  Problem  }

  \RemoveLine \RemoveHalfLine
  \begin{center}
    \begin{tikzpicture}[xscale=1.75,yscale=1.75]
      \node[dNode,label=above:\tiny{-10}] (l1) at (0.0,2.8) {\tikztext{L.NA}};
      \node[dNode,label=above:\tiny{-10}] (l2) at (0.0,2.1) {\tikztext{L.SA}};
      \node[dNode,label=above:\tiny{-10}] (l3) at (0.0,1.4) {\tikztext{L.EU}};
      \node[dNode,label=above:\tiny{-10}] (l4) at (0.0,0.7) {\tikztext{L.AS}};
      \node[dNode,label=above:\tiny{-10}] (l5) at (0.0,0.0) {\tikztext{L.AF}};

      \node[dNode,label=above:\tiny{-20}] (w1) at (3.0,2.8) {\tikztext{H.NA}};
      \node[dNode,label=above:\tiny{-20}] (w2) at (3.0,2.1) {\tikztext{H.SA}};
      \node[dNode,label=above:\tiny{-20}] (w3) at (3.0,1.4) {\tikztext{H.EU}};
      \node[dNode,label=above:\tiny{-20}] (w4) at (3.0,0.7) {\tikztext{H.AS}};
      \node[dNode,label=above:\tiny{-20}] (w5) at (3.0,0.0) {\tikztext{H.AF}};

      \node[node,label=above:\tiny{12}](m1) at (1.5,2.8) {\tikztext{NA}};
      \node[node,label=above:\tiny{8}] (m2) at (1.5,2.1) {\tikztext{SA}};
      \node[node,label=above:\tiny{14}](m3) at (1.5,1.4) {\tikztext{EU}};
      \node[node,label=above:\tiny{16}](m4) at (1.5,0.7) {\tikztext{AS}};
      \node[node,label=above:\tiny{7}] (m5) at (1.5,0.0) {\tikztext{AF}};

      \node (lName) at (0.0,-0.45) {\tikztext{L. Plant}};
      \node (wName) at (3.0,-0.45) {\tikztext{H. Plant}};
      \node (mName) at (1.5,-0.45) {\tikztext{Market}};

      \draw[eEdge] (w1) -- node[edgeText] {}  (m1) ;
      \draw[eEdge] (w1) -- node[edgeText] {}  (m2) ;
      \draw[eEdge] (w1) -- node[edgeText] {}  (m3) ;
      \draw[eEdge] (w1) -- node[edgeText] {}  (m4) ;
      \draw[eEdge] (w1) -- node[edgeText] {}  (m5) ;
      \draw[eEdge] (w5) -- node[edgeText] {}  (m1) ;
      \draw[eEdge] (w5) -- node[edgeText] {}  (m2) ;
      \draw[eEdge] (w5) -- node[edgeText] {}  (m3) ;
      \draw[eEdge] (w5) -- node[edgeText] {}  (m4) ;
      \draw[eEdge] (w5) -- node[edgeText] {}  (m5) ;
      \draw[eEdge] (l2) -- node[edgeText] {}  (m1) ;
      \draw[eEdge] (l2) -- node[edgeText] {}  (m2) ;
      \draw[eEdge] (l2) -- node[edgeText] {}  (m3) ;
      \draw[eEdge] (l2) -- node[edgeText] {}  (m4) ;
      \draw[eEdge] (l2) -- node[edgeText] {}  (m5) ;
      \draw[eEdge] (l4) -- node[edgeText] {}  (m1) ;
      \draw[eEdge] (l4) -- node[edgeText] {}  (m2) ;
      \draw[eEdge] (l4) -- node[edgeText] {}  (m3) ;
      \draw[eEdge] (l4) -- node[edgeText] {}  (m4) ;
      \draw[eEdge] (l4) -- node[edgeText] {}  (m5) ;
    \end{tikzpicture}
  \end{center}

  \note{
    \begin{itemize}
      \item a
    \end{itemize}
  }
\end{frame}

%=================================================================================================
\begin{frame}
  \frametitle{Notation of CPL  Problem }

  \textbf{Notation \& Decision variables}
  \begin{eqnarray*}
    \mathcal{I} &=& \mbox{set of potential plants (capability $\times$ location)}\\
    \mathcal{J} &=& \mbox{set of markets}\\
    D_j         &=& \mbox{demands at market $j$}\\
    K_i         &=& \mbox{capacities of plant $i$}\\
    f_i         &=& \mbox{construction costs of plant $i$}\\
    c_{ij}      &=& \mbox{cost of prod. and transp from plant $i$ to market $j$}\\
    y_{i}       &=& \mbox{plant construction binary decision variable}\\
                & & y_i = \begin{cases}
                          1~& \mbox{if plant $i$ is constructed}\\
                          0~& \mbox{otherwise}
                          \end{cases} \\
    x_{ij}      &=& \mbox{production transported between plant $i$ and market $j$} \\
  \end{eqnarray*}

  \note{
    \begin{itemize}
      \item what is the optimal course of action to fulfil demand
    \end{itemize}
  }
\end{frame}



%=================================================================================================
\begin{frame}
  \frametitle{Model of CPL Problem }

  \begin{alignat*}{4}
    \min\ z &= \sum_{i \in\ \mathcal{I}} f_i \, y_{i} + \sum_{i \in \mathcal{I}} \sum_{j \in \mathcal{J}} c_{ij} x_{ij} \\
    \mbox{s.t.,}\\
    \mbox{Supply $i$;~~~}& \sum_{j \in \mathcal{J}} x_{i,j} & \leq\ & K_i \, y_i & \qquad \forall\ i \in \mathcal{I} \\
    \mbox{Demand $j$;~~~}& \sum_{i \in \mathcal{I}} x_{i,j} ~~& =\ & D_j & \qquad \forall\ i \in \mathcal{I} \\
    \mbox{Non-Neg;~~~}& ~~~x_{i,j} &\geq\ & 0 & \qquad \forall\ (i,j) \in \mathcal{I} \times \mathcal{J} \\
    \mbox{Binary;~~~}&  ~~~y_{i} &\in\ & \{0,1\} & \qquad \forall\ i \in \mathcal{I} \\
  \end{alignat*}

  \note{
    \begin{itemize}
      \item \textbf{add 1:} limit \# H. Cap Plant or \# Low Cap Plant
      \item \textbf{add 2:} limit \# location
      \item \textbf{add 3:} condition market and location
    \end{itemize}
  }
\end{frame}

%========================================================================================================
\section[TSP]{Traveling Salesman Problem}
%========================================================================================================
\begin{frame}
  \frametitle{Intro to TSP }

  \RemoveLine \RemoveHalfLine
  \begin{center}
    \begin{tikzpicture}[xscale=1.5,yscale=1.5]
      \node[node] (n1) at (1.5,3.0)  {\tikztext{1}};
      \node[node] (n2) at (3.0,2.25) {\tikztext{2}};
      \node[node] (n3) at (3.0,0.75) {\tikztext{3}};
      \node[node] (n4) at (1.5,0.0)  {\tikztext{4}};
      \node[node] (n5) at (0.0,0.75) {\tikztext{5}};
      \node[node] (n6) at (0.0,2.25) {\tikztext{6}};
    \end{tikzpicture}
  \end{center}

  \begin{itemize}
      \item \textbf{Idea:} Given locations, find the shortest distances that all locations once
      \item \textbf{Important:} well-studied because fundamental of routing
  \end{itemize}

  \note{
    \begin{itemize}
      \item
    \end{itemize}
  }
\end{frame}

%========================================================================================================
\begin{frame}
  \frametitle{Natural Formulation \forIE}

  \small{
    \begin{eqnarray*}
      \mathcal{N} &=& \mbox{set of locations}\\
      \mathcal{E} &=& \mbox{set of connected locations (edges)}\\
      c_{i,j}     &=& \mbox{distances from location $i$ to location $j$}\\
      x_{i,j}     &=& \mbox{binary decision variable}\\
                  & & x_{i,j} = \begin{cases}
                                1~& \mbox{if edge ($i,j$) is used}\\
                                0~& \mbox{otherwise}\\
                                \end{cases}
    \end{eqnarray*}


    \begin{alignat*}{2}
    \min\ z     & =& \sum_{(i,j) \in\ \mathcal{E} } c_{i,j} x_{i,j}
    \end{alignat*}

    \begin{alignat*}{4}
      \mbox{s.t.} & & & \\
      \mbox{node n;~~} &  \sum_{i \in\ \mathcal{N} } x_{i,n} + \sum_{j \in\ \mathcal{N} } x_{n,j} ~ &=\ & 2 & \qquad \forall\ n \in\ \mathcal{N}\\
      \mbox{non-neg.;~~}&  ~~~~~x_{i,j} &\in\ & \{0,1\} & \qquad \forall\ (i,j) \in\ \mathcal{E}
    \end{alignat*}
  }

  \note{
    \begin{itemize}
      \item a
    \end{itemize}
  }
\end{frame}


%========================================================================================================
\begin{frame}
  \frametitle{Example of TSP }

  \begin{center}
    \begin{tikzpicture}[xscale=1.5,yscale=1.5]
      \node[node] (n1) at (1.5,3.0)  {\tikztext{1}};
      \node[node] (n2) at (3.0,2.25) {\tikztext{2}};
      \node[node] (n3) at (3.0,0.75) {\tikztext{3}};
      \node[node] (n4) at (1.5,0.0)  {\tikztext{4}};
      \node[node] (n5) at (0.0,0.75) {\tikztext{5}};
      \node[node] (n6) at (0.0,2.25) {\tikztext{6}};

      \only<1|handout:0>{
      \draw[eEdge] (n1) -- (n2) ;
      \path[eEdge] (n2) edge[bend left] (n4)    ;
      \draw[eEdge] (n4) -- (n5) ;
      \draw[eEdge] (n5) -- (n3) ;
      \draw[eEdge] (n3) -- (n6) ;
      \draw[eEdge] (n6) -- (n1) ;
      }

      \only<2|handout:0>{
      \draw[eEdge] (n1) -- (n2) ;
      \draw[eEdge] (n2) -- (n3) ;
      \draw[eEdge] (n3) -- (n1) ;
      \draw[eEdge] (n5) -- (n4) ;
      \draw[eEdge] (n4) -- (n6) ;
      \draw[eEdge] (n6) -- (n5) ;
      }

      \only<3|handout:0>{
      \draw[eEdge] (n1) -- (n2) ;
      \draw[eEdge] (n2) -- (n3) ;
      \draw[eEdge] (n3) -- (n4) ;
      \draw[eEdge] (n4) -- (n5) ;
      \draw[eEdge] (n5) -- (n6) ;
      \draw[eEdge] (n6) -- (n1) ;
      }
    \end{tikzpicture}
  \end{center}

  \begin{itemize}
    \item \textbf{What:} A 'difficult' ($\mathcal{NP}$-hard) combinatorial problem
    \item \textbf{Why:} number of \alert{sub-tour} ($2^{n-2}$)
    \item \textbf{How to solve:} taboo search, ant colony
  \end{itemize}

  \note{
    \begin{itemize}
      \item \textbf{What:} A 'difficult' ($\mathcal{NP}$-hard) combinatorial problem
      \item \textbf{Why:} number of \alert{sub-tour} ($2^{n-2}$)
      \item \textbf{How to solve:} taboo search, ant colony
    \end{itemize}
  }
\end{frame}


%========================================================================================================
\begin{frame}
  \frametitle{Sub-Tour Elimination \forIE}

  \begin{center}
    \begin{tikzpicture}[xscale=1.5,yscale=1.5]
      \node[node] (n1) at (1.5,3.0)  {\tikztext{1}};
      \node[node] (n2) at (3.0,2.25) {\tikztext{2}};
      \node[node] (n3) at (3.0,0.75) {\tikztext{3}};
      \node[node] (n4) at (1.5,0.0)  {\tikztext{4}};
      \node[node] (n5) at (0.0,0.75) {\tikztext{5}};
      \node[node] (n6) at (0.0,2.25) {\tikztext{6}};

      \draw[eEdge] (n1) -- (n2) ;
      \draw[eEdge] (n2) -- (n3) ;
      \draw[eEdge] (n3) -- (n1) ;
      \draw[eEdge] (n5) -- (n4) ;
      \draw[eEdge] (n4) -- (n6) ;
      \draw[eEdge] (n6) -- (n5) ;

    \end{tikzpicture}
  \end{center}

  \begin{itemize}
    \item \textbf{add more constraint} $x_{1,2}+x_{2,3}+x_{3,1} \leq\ 2$
  \end{itemize}

  \note{
    \begin{itemize}
      \item \textbf{What:} A 'difficult' ($\mathcal{NP}$-hard) combinatorial problem
      \item \textbf{Why:} number of \alert{sub-tour} ($2^{n-2}$)
      \item \textbf{How to solve:} taboo search, ant colony
    \end{itemize}
  }
\end{frame}

%========================================================================================================
\begin{frame}
  \frametitle{Practical Solution}

  \textbf{Observation}
  \begin{itemize}
    \item easy to find \alert{any TSP solution}
    \item difficult to find \alert{optimal solution} $\exists\ \frac{(n-1)!}{2}$ possible solution
    \item topological condition, particularly \alert{triangle inequality} reduce search space \NextLine
  \end{itemize}

  \textbf{2-Opt Heuristic}
  \begin{itemize}
    \item find any TSP solution (e.g., greedy algorithm)
    \item remove \alert{two edges} from the solution
    \item re-connect TSP with 'new' edges
  \end{itemize}

  \note{
    \begin{itemize}
      \item \textbf{What:} A 'difficult' ($\mathcal{NP}$-hard) combinatorial problem
      \item \textbf{Why:} number of \alert{sub-tour} ($2^{n-2}$)
      \item \textbf{How to solve:} taboo search, ant colony
    \end{itemize}
  }
\end{frame}

%==============================================================================================================
\section[VMI/VRP]{Vendor Managed Inventory and Vehicle Routing Problem}
%==============================================================================================================
  \begin{frame}
  \frametitle{Background of VMI}

  \begin{itemize}
    \item<1-> \textbf{What is Vendor Managed Inventory (VMI)?}
    \begin{itemize}
      \item<2-> Supplier orders/controls inventory for retailers
      \item<2-> A popular \alert{continuous replenishment} in retailer business
      \item<2-> Inventory is managed by vendor, not retailer \NextLine
    \end{itemize}
    \item<1-> \textbf{Implementation}
    \begin{itemize}
      \item<3-> \alert{Strong relationships}
      \item<3-> Information Technology $\rightarrow$ Coordination
      \item<3-> \alert{``Everyday'' products} \NextLine
    \end{itemize}
    \item<1-> \textbf{Benefits}
    \begin{itemize}
      \item<4-> \textbf{Both}: Achieving \alert{quick response} and more trust
      \item<4-> \textbf{Retailer}: Reducing labors and \alert{inventory} (space and capital)
      \item<4-> \textbf{Supplier}: Gaining competitive edge and LTL shipments \NextLine
    \end{itemize}
    \item<1-> \textbf{Example}
    \begin{itemize}
      \item<4-> Procter \&\ Gamble and WalMart, HP~, Johnson \&\ Johnson
    \end{itemize}
  \end{itemize}

  \note{
    \begin{itemize}
      \item \textbf{continuous replenishment}: not determined in advance
      \item Retailer-Supplier Logistics are virtually connected
      \item Retailer must have economy of scales and significance
      \item Retailer outsources purchasing to vendor (save purchasing)
      \item Supplier must consolidate order and do LTL shipping
      \item \textbf{dilemma}: less inventory for retailer, more inventory for vendor
      \item Dell pull inventory example
      \item \textbf{products candidate}: everyday use, non-perishable
    \end{itemize}
  }

  \end{frame}

%==============================================================================================================
  \begin{frame}
  \frametitle{What are operation issues behind VMI?}

  \begin{itemize}
    \item<1-> \textbf{VMI at face value}
    \begin{itemize}
      \item<2-> More information $\rightarrow$ real time delivery
      \item<2-> Seamless products transfer $\rightarrow$ eliminate receiving and inspecting \NextLine
    \end{itemize}
    \item<1-> \textbf{How to improve response time and reduce inventory?}
    \begin{itemize}
      \item<3-> \alert{Increasing frequency} of shipping
      \item<3-> \alert{Reducing} shipping \alert{quantities}
      \item<3-> Shipping to \alert{multiple retailers} \NextLine
    \end{itemize}
    \item<1-> \textbf{Operation issue}
    \begin{itemize}
      \item<4-> VMI $\rightarrow$ Vehicle Routing Problem (VRP)
    \end{itemize}
  \end{itemize}

  \note{
    \begin{itemize}
      \item \textbf{IT}: GPS, RFID, EDI
      \item \textbf{Trust \& Quality}: skin, hide, no quantity-per-lot discounted
      \item \textbf{At face value}: VMI doesn't help operation much, logistic operation
    \end{itemize}
  }

  \end{frame}

%========================================================================================================
\begin{frame}
  \frametitle{Vehicle Routing Problem \forISE}

  \RemoveLine
  \begin{center}
        \includegraphics[width=0.55\linewidth]{\figureLocation{scm_trans/exampleVRP.png}}
  \end{center}
  \begin{itemize}
    \item \textbf{What:} model describe routing of vehicle (from/to depot)
    \item \textbf{Idea:} minimizing total times/distances of multiple vehicles
  \end{itemize}

  \note{
    \begin{itemize}
      \item \textbf{What:} model describe the fastest \alert{path} to transport products from origin to destination
      \item \textbf{Idea:} minimizing total times/distances
    \end{itemize}
  }
\end{frame}

%==============================================================================================================
  \begin{frame}
  \frametitle{Background of VRP  \forIE}

  \begin{itemize}
    \item<1-> \textbf{Classification of Routing Problem}
    \begin{itemize}
      \item<2-> \textbf{Infinite capacity}: Traveling Salesman Problem (TSP)
      \item<2-> \textbf{Finite capacity}: TSP + Bin Packing Problem (BPP)
      \item<2-> \textbf{Time windows}: specific pick-up and drop-off times in each node
      \item<2->[] \hskip30pt TSP + BPP + Scheduling Problem
      \item<2-> \textbf{Stochastic}: demands or time windows occurred with some probability \NextLine
    \end{itemize}
    \item<1-> \textbf{Solution Method}
    \begin{itemize}
      \item<3-> \textbf{Exact solution:} Mixed Integer Programming
      \item<3-> \textbf{Heuristic:} \alert{Traditional} (e.g., 2-Opt, greedy) and \alert{Meta-Heuristic}
      \item<3-> \textbf{Approximated algorithm:}: Heuristic with known error/upperbound
    \end{itemize}
  \end{itemize}

  \note{
    \begin{itemize}
      \item Visit each and every node only one time $+$ sub-tour elimination ($2^{n-2}$)
      \item TSP is NP-hard; meaning difficult to solve
      \item TSP is important problem; many heuristic for start-up and improving step
      \item Bin-packing Problem is to pack/contain all different size item with the same size bin/bastket as less as possible
      \item Scheduling problem: Jobshop
    \end{itemize}
  }


  \end{frame}

%==============================================================================================================
  \begin{frame}
  \frametitle{Simple Coordination with VMI/VRP  } \RemoveLine
  \only<1|handout:0>{
    \begin{center}
      \includegraphics[width=0.85\textwidth]{\figureLocation{scm_trans/IdeaVMI01.png}}\\
    \end{center}
  }
  \only<2|handout:0>{
    \begin{center}
      \includegraphics[width=0.85\linewidth]{\figureLocation{scm_trans/IdeaVMI02.png}}\\
      \footnotesize{Cluster of retailers based on \alert{distance}}
    \end{center}
  }
  \only<3|handout:0>{
    \begin{center}
      \includegraphics[width=0.85\linewidth]{\figureLocation{scm_trans/IdeaVMI03.png}}\\
    \end{center}
  }
  \only<4|handout:1>{
    \begin{center}
      \includegraphics[width=0.85\linewidth]{\figureLocation{scm_trans/IdeaVMI04.png}}\\
      \footnotesize{Cluster of retailer based on \alert{inventory}}
    \end{center}
  }
  \note{
    \begin{itemize}
      \item What is vendor facing is Routing problem
    \end{itemize}
  }
  \end{frame}

%==============================================================================================================
  \begin{frame}
  \frametitle{VMI Potential Problems}

    \begin{itemize}
      \item \textbf{Pilot VMI:} SKUs, vendor selection, hi-involvement
      \item \textbf{Story behind data:} truckload, special order, demand pattern
      \item \textbf{Efficient communication:} new product, seasonal variability
      \item \textbf{Change management:} promotion, compensation
    \end{itemize}

    \NextHalfLine
    \textbf{VMI Success Factors}
    \begin{itemize}
      \item Focus your \alert{efforts} $\rightarrow$ competency \& accuracy
      \item Trust between vendor/retailer $\rightarrow$ willingness  \& patience
      \item Effective computer systems $\rightarrow$ \alert{communication}
    \end{itemize}


  \note{
    \begin{itemize}
      \item K-Mart and Whiehall-Robbin (Advile)
    \end{itemize}
  }
  \end{frame}


%================================================================================================-
\section[FacSelection]{Introduction of Facility Selection}
%================================================================================================-

\begin{frame}
  \frametitle{Importance of Sites and Key Questions}

  \begin{center}
    \begin{minipage}{0.75\linewidth}
    \begin{quotation}
      There are three things that matter in property: location, location, location.
    \end{quotation}
    \end{minipage}
  \end{center}

  \NextHalfLine
  \begin{itemize}
      \item Providing structure to the network
      \item Affecting inventory and \alert{transportation costs}
      \item Improving customer services
  \end{itemize}

  \NextLine
  \textsc{Key Questions}
  \begin{itemize}
    \item \alert{How many} facilities should there be? \uncover<2->{\hskip3ex Service \textsc{VS} Scale}
    \item \alert{Where} should they be located? \uncover<2->{\hskip3ex $\rightarrow$ copy, compute, wait-and-see}
    \item What \alert{size/cap} should they be? \uncover<2->{\hskip3ex $\rightarrow$ usage + expansion + buffer}
  \end{itemize}


  \note{
    \begin{itemize}
      \item \textbf{Economy}: cheaper to produce/order in large batches; setup cost and time
      \item \textbf{Uncertainty}: demand; lead times; delivery of supply
      \item \textbf{Speculation}: FG held in anticipation of a rise in their value or costs
      \item \textbf{Smoothing}: irregularities of demand pattern
      \item \textbf{Transportation and Logistics:} inventory in transit from one location to other
    \end{itemize}

    \begin{itemize}
    \item Managing capacity
    \begin{itemize}
      \item time flexibility from workforce (OT, pay per piece)
      \item seasonal workforce or sub-contracting
      \item designing product flexibility into production processes \NextHalfLine
    \end{itemize}
    \item Managing inventory
    \begin{itemize}
      \item common components across multiple products
      \item building inventory of high demand or predictable demand products
    \end{itemize}
  \end{itemize}

  }
\end{frame}

%============================================================================================================--
\begin{frame}
  \frametitle{Nature of Location Analysis}

  \begin{itemize}
    \item \textbf{Type of Facility}: cheaper to produce/order in large batches; setup cost and time. (man, activities, surroundings)
    \begin{itemize}
      \item \textsl{Plant:} driven by \alert{economics}
      \item \textsl{Retail:} driven by revenue.
      \item \textsl{Service:} driven by service factors.  \NextHalfLine
    \end{itemize}
    \item \textbf{Cost Structure:} investments (tech/cap)+ periodic upkeep (time/cap) + shipments (qty)
    \item \textbf{Clustering:} because of utility, market, infrastructure (motor terminals, shops)
    \item \textbf{Rent VS Own:} time horizon, land, location,
    \item \textbf{Present VS Future:} not average, nor max  \uncover<2->{\hskip3ex $\rightarrow$ overflow qty/ inventory}
  \end{itemize}

  \note{
    \begin{itemize}
      \item \textbf{Example:} fashion in ITY, H. auto in GRE, mutual fun in Boston
      \item \textbf{Example:} fashion in ITY, H. auto in GRE, mutual fun in Boston
      \item \textbf{Cluster:} concentration of \alert{interconnected} companies and institutions
      \item \textbf{Cluster \& Competition:} infrastructure, regulation
    \end{itemize}

    \begin{itemize}
    \item Managing capacity
    \begin{itemize}
      \item time flexibility from workforce (OT, pay per piece)
      \item seasonal workforce or sub-contracting
      \item designing product flexibility into production processes \NextHalfLine
    \end{itemize}
    \item Managing inventory
    \begin{itemize}
      \item common components across multiple products
      \item building inventory of high demand or predictable demand products
    \end{itemize}
  \end{itemize}

  }
\end{frame}

%============================================================================================================--
\begin{frame}
  \frametitle{Factors that affects site selection}

  \begin{itemize}
    \item \textbf{Objective:} min cost, max distance (nuclear plant), max coverage (police)
    \item \textbf{Demands and Supplies:} big cities, community, cluster of 'things' (CDC, Gift)
    \item \textbf{Distribution and Growth:} traffic type (Bus Stop VS BTS station)
    \item \textbf{Infrastructure:} road, highway, water, electricity, data
    \item \textbf{Other:} land price, potential,
    \item \textbf{Regulation \& Benefit:} BOI, zoning
    \begin{itemize}
      \item \textbf{\code{Must-Have:}} water treatment, waterway, electricity
      \item \textbf{\code{Must-Avoid:}} restricted/religious/historical areas, community
      \item \textbf{\code{Zone:}} industrial park, land use zoning, flood plane, buffering
      \item \textbf{\code{Civic Codes:}} EIA, EHIA, license
    \end{itemize}
  \end{itemize}

  \note{
    \begin{itemize}
    \item \textbf{Rectilinear Distance:} Manhattan Grid
    \item \textbf{Straight Line Distance:}
    \item \textbf{Chebyshev Distance:}
    \item \textbf{Actual Distance:}
    \end{itemize}
  }
\end{frame}


%============================================================================================================--
%\begin{frame}
%  \frametitle{Distance Measurement}
%  \RemoveHalfLine
%  \begin{center}   \scriptsize{
%  \begin{tabular}{ccc}
%      \includegraphics[width=0.225\textwidth]{\figureLocation{scm_loc/dist01.png}}
%    &
%      \includegraphics[width=0.225\textwidth]{\figureLocation{scm_loc/dist02.png}}
%    &
%      \includegraphics[width=0.225\textwidth]{\figureLocation{scm_loc/dist03.png}}
%    \\
%      \textbf{Rectilinear} (reach truck in aisle)
%    &
%      \textbf{Straight Line} (forklift truck in crossdock)
%    &
%      \textbf{Chebyshev} (AS/RS)
%    \\
%      \includegraphics[width=0.225\textwidth]{\figureLocation{scm_loc/dist05.png}}
%    &
%      \begin{minipage}{0.25\textwidth}
%        \RemoveLine \RemoveLine \RemoveLine
%        \textsc{\code{Property}}\\
%        \begin{eqnarray*}
%          d: \mathbf{M} \times \mathbf{M} &\rightarrow & \mathbb{R} \\
%          d(\mathbf{x},\mathbf{x}) &=& 0 \\
%          d(\mathbf{x},\mathbf{y}) &=& d(\mathbf{y},\mathbf{x}) \\
%          d(\mathbf{x},\mathbf{z}) &\leq& d(\mathbf{x},\mathbf{y}) +d(\mathbf{y},\mathbf{z}) \\
%        \end{eqnarray*}
%      \end{minipage}
%    &
%      \includegraphics[width=0.225\textwidth]{\figureLocation{scm_loc/dist04.png}} \\
%    \\
%
%      \textbf{Great Circle} (country distance)
%    &
%
%    &
%      \textbf{Actual} (city road distance) \\
%
%  \end{tabular} }
%  \end{center}
%
%  \note{
%    \begin{itemize}
%    \item \textbf{Rectilinear Distance:} Manhattan Grid
%    \item \textbf{Straight Line Distance:}
%    \item \textbf{Chebyshev Distance:}
%    \item \textbf{Actual Distance:}
%    \end{itemize}
%  }
%\end{frame}

%============================================================================================================-
\begin{frame}
  \frametitle{Questions for facility location}

  \begin{itemize}
    \item How many facilities should there be?
    \item Where should they be located?
    \item How does each facility function together?
  \end{itemize}

  \NextHalfLine
  \begin{alertblock}{}{}
    Theoretically, the question is strategic, \alert{multi-facet}, and difficult to optimize \uncover<2->{ $\rightarrow$ MILP, AHP, fuzzy logic }
  \end{alertblock}

  \NextLine
  \textsc{Practical Good news}
  \begin{itemize}
    \item \textbf{Cluster:} city, population, airport, utility, workforce
    \item \textbf{Competitors:} e.g., Boots \textbf{VS} Watson, 7Eleven \textbf{VS} 'other'
    \item \textbf{Wait-And-See approach:} one warehouse/one building at a time
    \item \textbf{Issues:} expanding urban area, restricted zone
  \end{itemize}

  \note{
    \textbf{Purpose:} overview facility location problem without calculation   \\
    \textbf{Leading Question:} What are criteria of good warehouse? \\
    \begin{itemize}
      \item \textbf{how many:} wait-and-see one at time
      \item \textbf{model:} tool, not definitive guide (disaster, disruption, vision )
      \item \textbf{size:} never enough space $\rightarrow$ inventory control
      \item \textbf{customer order:} de-central $\rightarrow$ store only some items; customer order $\rightarrow$ items located;
      \item \textbf{problems:} what should each warehouse
    \end{itemize}
  }
  \end{frame}

%================================================================================================-

\begin{frame}
  \frametitle{Reference}

  \begin{footnotesize}
    \bibliography{\refLocation{Coordination}}
  \end{footnotesize}

  \note{
    [empty] note for citation page
  }

\end{frame}

%================================================================================================-
\end{document}


%==============================================================================================================
  \begin{frame}
  \frametitle{Benefits of VMI }

  \begin{center}
  \small{
  \begin{tabular}{ll}\hline
    \textbf{Benefits to Vendors}  & \textbf{Benefit to Retailer} \\ \hline
    $\bullet$  lower Inventory investment    & $\bullet$ higher inventory turnover  \\
    (smaller batch)               & (higher shipping frequency) \\
    $\bullet$ better scheduling \& planning & $\bullet$ better market information \\
    $\bullet$ better market information     & $\bullet$ more optimal product mixes \\
    $\bullet$ close customer ties           & $\bullet$ less replenishing cost \\ \hline
  \end{tabular}
  }
  \end{center}

  \NextHalfLine
  \textbf{Success Stories}
  \begin{itemize}
      \item<2-> Customer service at K-mart improves from $80+$ to $90+$ and inventory turnover improves from 3(12) to 10(17)
      \item<2-> ACE Hardware has seen fill rates rise 4\% to 96\% in past few years
      \item<2-> Oshawa Food involves Pillsbury, Quaker and Heniz in VMI and improves its customer service level 99/%
  \end{itemize}

  \note{
    \begin{itemize}
      \item ACE
      \item \textbf{Trust \& Quality}: skin, hide, no quantity-per-lot discounted
      \item \textbf{At face value}: VMI doesn't help operation much, logistic operation
    \end{itemize}
  }
  \end{frame}
